% Vietnamese translation for OI Public Library
% Source-File: IOI中国国家候选队论文/2006/俞鑫/俞鑫.ppt
\documentclass[11pt]{article}
\usepackage{oipl-vn}

\title{Bàn cờ trong bàn cờ\\{\large Ghi chú theo slide}}
\author{Yu Xin}
\date{2006}

\begin{document}
\maketitle

\origtitle{Chessboard partition thinking}
\sourcepath{IOI China candidate team papers / 2006 / Yu Xin / Yu Xin.ppt}

\section{Phạm vi}

Bộ slide đi cùng bài viết về tư tưởng phân hoạch bàn cờ. Bản ghi này dựa trên
bản bài viết đã dịch, giữ lại đường dây trình bày chính: mô hình bàn cờ, cách
chia bàn thành các vùng con, và hai ví dụ tiêu biểu.

\section{Phân hoạch bàn cờ}

Một bàn cờ có thể được chia theo hàng, cột, góc phần tư, đường chéo, màu ô,
hoặc theo những vùng được xác định bởi trạng thái bài toán. Điểm quan trọng
không phải là chia nhỏ tùy ý, mà là chia sao cho mỗi phần con cùng dạng với
bài toán ban đầu hoặc làm lộ ra một bất biến dễ kiểm soát.

\section{Phủ bàn cờ}

Ví dụ đầu là bài phủ bàn \(2^k\times 2^k\) có một ô đặc biệt bằng các quân
chữ L. Chia bàn thành bốn góc phần tư. Góc phần tư chứa ô đặc biệt giữ nguyên
vai trò ấy; ba góc còn lại được tạo một ô đặc biệt giả bằng cách đặt một quân
chữ L ở vùng trung tâm. Sau đó bốn bài con độc lập và có cùng cấu trúc với
bài gốc.

Số quân cần dùng là
\[
  \frac{4^k-1}{3},
\]
vì mỗi quân phủ ba ô và đúng một ô của bàn không cần phủ. Đây là ví dụ điển
hình cho tư tưởng chia để trị trên bàn cờ.

\section{Cờ Khổng Minh}

Ví dụ thứ hai phức tạp hơn: không chỉ chia theo kích thước, ta còn phải nhận
ra trạng thái nào của từng vùng đủ để ghép nghiệm toàn cục. Bài viết dùng mô
hình cờ Khổng Minh để nhấn mạnh rằng một phép chia tốt thường đi kèm quy tắc
ghép ngược lại. Nếu mỗi vùng con được giải độc lập nhưng không có thông tin
biên phù hợp, lời giải cục bộ sẽ không tạo được lời giải toàn cục.

\section{Ghi nhớ}

Khi gặp bài toán bàn cờ, nên thử tìm một cách chia làm giảm kích thước mà vẫn
bảo toàn cấu trúc. Các ô đặc biệt, màu ô, vùng biên và quan hệ đối xứng là
những tín hiệu thường quyết định phép phân hoạch đúng.

\end{document}
